\FigureLayoutDeclare{img/1977/shandong_b/q10_solution_fig1.png}{7.5cm}{36bp 0bp 13bp 111bp}
\FigureLayoutDeclare{img/1977/shandong_b/q6_fig1.png}{7.0cm}{41bp 36bp 88bp 39bp}
\FigureLayoutDeclare{img/1977/shandong_b/q7_solution_fig1.png}{9.0cm}{0bp 68bp 3bp 0bp}

\chapter{1977年山东B卷}

\section{解答题}

\begin{problem}
\(\sqrt{\left( -a \right)^2}=?\)
\end{problem}

\begin{solution}
因为 \(\left( -a \right)^2=a^2\)，而算术平方根的值非负，所以
\(\sqrt{\left( -a \right)^2}=\sqrt{a^2}=\lvert a\rvert\). 当 \(a\geqslant0\) 时，\(\lvert a\rvert=a\)；当 \(a<0\) 时，\(\lvert a\rvert=-a\). 因此
\[
\sqrt{\left( -a \right)^2}=\begin{cases}
a, & a\geqslant0,\\
-a, & a<0.
\end{cases}
\]
\end{solution}

\begin{problem}
\(\log_{\frac{1}{2}}\left[ \left( 1-\cos^2\frac{\pi}{6} \right)\left( \cot\frac{\pi}{4} \right)^{-2} \right]=?\)
\end{problem}

\begin{solution}
由 \(\cos\frac{\pi}{6}=\frac{\sqrt{3}}{2}\) 和 \(\cot\frac{\pi}{4}=1\)，得
\(1-\cos^2\frac{\pi}{6}=1-\frac{3}{4}=\frac{1}{4}\)，\(\left( \cot\frac{\pi}{4} \right)^{-2}=1^{-2}=1\).
所以原式 \(=\log_{\frac{1}{2}}\frac{1}{4}\). 又因为 \(\frac{1}{4}=\left(\frac{1}{2}\right)^2\)，故原式 \(=2\).
\end{solution}

\begin{problem}
写出下列函数的定义域：

\begin{enumerate}
\item \(y=\frac{\sqrt{\lvert x-3\rvert}}{x-2}\)；
\item \(y=\lg\left( x+3 \right)\).
\end{enumerate}
\end{problem}

\begin{solution}
\begin{enumerate}
\item 对任意实数 \(x\)，都有 \(\lvert x-3\rvert\geqslant0\)，所以根式始终有意义；分母不能为零，需满足 \(x-2\ne0\). 因此定义域为 \(\{x\in\R\mid x\ne2\}\).

\item 对数的真数必须为正，需满足 \(x+3>0\)，所以定义域为 \(\{x\in\R\mid x>-3\}\).
\end{enumerate}
\end{solution}

\begin{problem}
写出过 \((0,1)\) 和 \((2,3)\) 两点的直线方程.
\end{problem}

\begin{solution}
两点的横坐标不同，直线斜率为
\[
k=\frac{3-1}{2-0}=1
\]
由点斜式，所求直线满足 \(y-1=1\left( x-0 \right)\)，即
\(y=x+1\)，或写成一般式 \(x-y+1=0\).
\end{solution}

\begin{problem}
写出双曲线 \(\frac{x^2}{16}-\frac{y^2}{9}=1\) 的渐近线方程.
\end{problem}

\begin{solution}
将方程写成标准形式
\[
\frac{x^2}{4^2}-\frac{y^2}{3^2}=1
\]
因此双曲线的参数为 \(a=4\)，\(b=3\). 对于双曲线 \(\frac{x^2}{a^2}-\frac{y^2}{b^2}=1\)，其渐近线为 \(y=\pm\frac{b}{a}x\). 代入 \(a\)，\(b\) 得
\[
y=\pm\frac{3}{4}x
\]
即渐近线方程为 \(\frac{x}{4}+\frac{y}{3}=0\) 和 \(\frac{x}{4}-\frac{y}{3}=0\).
\end{solution}

\begin{problem}
如图所示的二视图是什么几何体？

\bitmapfigure[float=inline,width=6.0cm]{img/1977/shandong_b/q6_fig1.png}
\end{problem}

\begin{solution}
图中的一个视图为正方形，另一个视图为以一个顶点为顶点、以正方形的一条边为底边的等腰三角形. 这说明几何体的底面是正方形，四个侧面都是以同一个顶点为公共顶点的三角形，且顶点在底面中心的正上方. 因此该几何体是正四棱锥.
\end{solution}

\begin{problem}
平行四边形 \(ABCD\) 的两邻边之比为 \(1:2\)，\(M\) 为大边 \(AB\) 的中点，连接 \(MD\)，\(MC\)，求证：\(\angle DMC=\frac{\pi}{2}\).
\end{problem}

\begin{solution}
由题意，\(AB:AD=2:1\)，且 \(M\) 是 \(AB\) 的中点，所以
\[
AM=MB=\frac{1}{2}AB=AD
\]
在等腰三角形 \(ADM\) 中，\(AM=AD\)，故
\(\angle ADM=\angle AMD\). 又因为 \(AB\parallel DC\)，且 \(A\)，\(M\)，\(B\) 三点共线，所以直线 \(AM\parallel DC\)，从而由平行线所成的角相等，\(\angle AMD=\angle MDC\). 因此
\[
\angle ADC=\angle ADM+\angle MDC=2\angle MDC
\]

\solutionfigure{\bitmapfigure[width=6.0cm]{img/1977/shandong_b/q7_solution_fig1.png}}

平行四边形的对边相等，故 \(AD=BC\)，于是 \(MB=BC\). 在等腰三角形 \(MBC\) 中，\(MB=BC\)，所以 \(\angle BMC=\angle MCB\). 由于 \(MB\parallel DC\)，有 \(\angle BMC=\angle MCD\)，从而
\[
\angle BCD=\angle MCB+\angle MCD=2\angle MCD
\]

平行四边形的邻角互补，即 \(\angle ADC+\angle BCD=\pi\). 因此
\[
\angle MDC+\angle MCD
=\frac{1}{2}\left( \angle ADC+\angle BCD \right)
=\frac{\pi}{2}
\]
在三角形 \(DMC\) 中，三角形内角和为 \(\pi\)，所以
\[
\angle DMC=\pi-\angle MDC-\angle MCD=\frac{\pi}{2}
\]
\end{solution}

\begin{problem}
\(\frac{\sin\alpha}{2}x^2+x+\frac{1}{2\cos\alpha}=0\)（\(0<\alpha<\frac{\pi}{2}\)）为关于 \(x\) 的二次方程.

\begin{enumerate}
\item 当 \(\alpha\) 为何值时，方程有两个不相等的实数根？
\item 它的根是什么？
\end{enumerate}
\end{problem}

\begin{solution}
\begin{enumerate}
\item 因为 \(0<\alpha<\frac{\pi}{2}\)，所以 \(\sin\alpha>0\)，\(\cos\alpha>0\)，该方程确为关于 \(x\) 的二次方程. 其二次项系数、一次项系数和常数项分别为
\(A=\frac{\sin\alpha}{2}\)，\(B=1\)，\(C=\frac{1}{2\cos\alpha}\).
方程有两个不相等的实数根的充要条件是判别式 \(\Delta=B^2-4AC>0\). 计算得
\[
\Delta=1-4\cdot\frac{\sin\alpha}{2}\cdot\frac{1}{2\cos\alpha}
=1-\frac{\sin\alpha}{\cos\alpha}=1-\tan\alpha
\]
所以需要 \(\tan\alpha<1\). 在区间 \(\left(0,\frac{\pi}{2}\right)\) 内，正切函数递增，且 \(\tan\frac{\pi}{4}=1\)，故
\[
0<\alpha<\frac{\pi}{4}
\]

\item 当 \(0<\alpha<\frac{\pi}{4}\) 时，\(\Delta=1-\tan\alpha>0\). 由求根公式
\[
x=\frac{-B\pm\sqrt{\Delta}}{2A}
=\frac{-1\pm\sqrt{1-\tan\alpha}}{\sin\alpha}
\]
因此两个不相等的实数根为
\(x_1=\frac{-1+\sqrt{1-\tan\alpha}}{\sin\alpha}\)，\(x_2=\frac{-1-\sqrt{1-\tan\alpha}}{\sin\alpha}\).
\end{enumerate}
\end{solution}

\begin{problem}
试比较等底等高的正四棱锥和正圆锥的侧面积，看哪个的侧面积较小.
\end{problem}

\begin{solution}
设两几何体的底面积均为 \(S>0\)，高均为 \(h>0\). 对正四棱锥，设正方形底面的边长为 \(a\)，则 \(a^2=S\)，从顶点到底面某一边中点的斜高为
\[
l_1=\sqrt{h^2+\left( \frac{a}{2} \right)^2}=\sqrt{h^2+\frac{S}{4}}
\]
因此正四棱锥的侧面积为
\[
S_1=4\cdot\frac{1}{2}a l_1
=2\sqrt{S}\sqrt{h^2+\frac{S}{4}}
\]
从而
\[
S_1^2=S^2+4Sh^2
\]

对正圆锥，设底面半径为 \(r\)，母线长为 \(l_2\). 由 \(\pi r^2=S\) 和勾股定理，得
\(r=\sqrt{\frac{S}{\pi}}\)，\(l_2=\sqrt{h^2+r^2}=\sqrt{h^2+\frac{S}{\pi}}\).
所以正圆锥的侧面积为
\[
S_2=\pi r l_2
\]
从而
\[
S_2^2=\pi^2\cdot\frac{S}{\pi}\left( h^2+\frac{S}{\pi} \right)=S^2+\pi Sh^2
\]

因为 \(\pi<4\) 且 \(S\)，\(h>0\)，所以 \(S_2^2<S_1^2\). 两个侧面积均为正数，故 \(S_2<S_1\). 因此正圆锥的侧面积较小.
\end{solution}

\begin{problem}
一只船以 \(v\text{ 海里/小时}\) 的速度向正北航行，在 \(A\) 处望见灯塔在船的北 \(\alpha\) 度东，继续航行 \(t\text{ 小时}\) 后，在 \(B\) 处望见 \(S\) 在船的北 \(\beta\) 度东，求第二次望见灯塔 \(S\) 时船和灯塔的距离 \(BS\).
\end{problem}

\begin{solution}
\solutionfigure{\bitmapfigure[width=4.2cm]{img/1977/shandong_b/q10_solution_fig1.png}}

题面中的 \(\alpha\)、\(\beta\) 按度数表示.
船从 \(A\) 到 \(B\) 向正北航行，所以
\[
AB=vt
\]
在三角形 \(ABS\) 中，\(AB\) 的方向是正北，故 \(\angle BAS=\alpha\). 在 \(B\) 点，射线 \(BA\) 指向正南，而射线 \(BS\) 是北 \(\beta\) 度东，所以
\[
\angle ABS=180^\circ-\beta
\]
由三角形内角和，
\[
\angle ASB=180^\circ-\angle BAS-\angle ABS
=180^\circ-\alpha-\left( 180^\circ-\beta \right)=\beta-\alpha
\]
由图中灯塔位于航线东侧且在 \(B\) 的北方可知 \(\beta>\alpha\)，因此该角为正. 在三角形 \(ABS\) 中应用正弦定理，\(BS\) 对应角 \(\alpha\)，\(AB\) 对应角 \(\beta-\alpha\)，于是
\[
\frac{BS}{\sin\alpha}=\frac{AB}{\sin\left( \beta-\alpha \right)}
=\frac{vt}{\sin\left( \beta-\alpha \right)}
\]
所以
\[
BS=\frac{vt\sin\alpha}{\sin\left( \beta-\alpha \right)}
\]
\end{solution}

\begin{problem}
某隧道的横断面如图所示，已知隧道横断面的周长为 \(23\text{ 米}\)，试根据图所示尺寸，求：

\begin{enumerate}
\item 当 \(x\) 和 \(y\) 为何值时，隧道的横断面面积最大？
\item 最大面积是多少？
\end{enumerate}

（取 \(\pi\approx3\)）（图尺寸单位：米）

\bitmapfigure[width=4.8cm]{img/1977/shandong_b/q11_fig1.png}
\end{problem}

\begin{solution}
\begin{enumerate}
\item 设半圆半径为 \(x\) 米，两侧直墙的高度为 \(y\) 米，横断面面积为 \(S\) 平方米. 按图中尺寸，半圆弧长取 \(\pi x\approx3x\). 底部边界由两段长度为 \(0.5\) 米的水平线、凹入部分的两条竖直边（每条长 \(0.5\) 米）以及凹入部分的上边组成；其水平边界总长为直径 \(2x\)，竖直边总长为 \(0.5\times2\). 因此周长条件给出
\[
23\approx3x+2y+2x+0.5\times2
\]
从而
\[
y\approx\frac{23-3x-2x-0.5\times2}{2}=11-2.5x
\]

隧道横断面面积等于上半圆面积加上两侧墙之间的矩形面积，再减去底部凹入的矩形面积. 其中凹入部分的宽为 \(2x-1\) 米，高为 \(0.5\) 米，所以
\[
S\approx\frac{3x^2}{2}+2xy-\left( 2x-1 \right)\times0.5
\]
代入 \(y=11-2.5x\)，化简得
\[
S\approx\frac{3x^2}{2}+2x\left( 11-2.5x \right)-\left( 2x-1 \right)\times0.5=-3.5x^2+21x+0.5
\]
几何尺寸还需满足 \(2x-1>0\) 且 \(y>0\)，即 \(\frac{1}{2}<x<\frac{22}{5}\). 在这个范围内，配方得
\[
S\approx-3.5\left( x-3 \right)^2+32
\]
因为 \(-3.5\left( x-3 \right)^2\leqslant0\)，且 \(x=3\) 在允许范围内，所以当 \(x=3\) 米时，面积取得最大值；此时
\[
y=11-2.5\times3=3.5\text{ 米}
\]

\item 由上式直接得到
\[
S_{\max}=32\text{ 平方米}
\]
\end{enumerate}
\end{solution}

\section{参考题}

\begin{problem}
求极限：\(\lim\limits_{x\to0}\frac{\ln\left( 1+x \right)}{x}\).
\end{problem}

\begin{solution}
当 \(x\) 趋近于 \(0\) 时，\(1+x>0\)，所以对数函数在所需邻域内有定义. 此时分子 \(\ln\left( 1+x \right)\to0\)，分母 \(x\to0\)，且分子、分母在该邻域内可导. 由洛必达法则，
\[
\lim\limits_{x\to0}\frac{\ln\left( 1+x \right)}{x}
=\lim\limits_{x\to0}\frac{\dfrac{1}{1+x}}{1}
=\lim\limits_{x\to0}\frac{1}{1+x}=1
\]
\end{solution}

\begin{problem}
用定积分求抛物线 \(y^2=4x\) 与 \(x^2=4y\) 所围成的曲边形面积.
\end{problem}

\begin{solution}
\solutionfigure{\bitmapfigure[width=5.0cm]{img/1977/shandong_b/ref1_q2_solution_fig1.png}}

由 \(y^2=4x\) 和 \(x^2=4y\) 可知交点都在第一象限或原点. 联立两式，将 \(y=\frac{x^2}{4}\) 代入 \(y^2=4x\)，得
\[
\frac{x^4}{16}=4x
\quad\Longleftrightarrow\quad
x\left( x^3-64 \right)=0
\]
所以交点的横坐标为 \(x=0\) 和 \(x=4\). 在 \(0\leqslant x\leqslant4\) 上，上方曲线为 \(y=2\sqrt{x}\)，下方曲线为 \(y=\frac{x^2}{4}\). 因此所围成区域的面积为
\[
S=\int_0^4\left( 2\sqrt{x}-\frac{x^2}{4} \right)\,\mathrm{d}x
=\left.\frac{4}{3}x^{\frac{3}{2}}\right|_0^4-\left.\frac{1}{12}x^3\right|_0^4
=\frac{32}{3}-\frac{16}{3}=\frac{16}{3}
\]
所以曲边形面积为 \(\frac{16}{3}\) 平方单位.
\end{solution}
